\section{Literature Review}

\subsection{Scope}

This project analyses the \emph{Olist Brazilian E-Commerce} dataset: nine
relational tables covering 99{,}441 orders placed on a Brazilian marketplace
between September 2016 and October 2018 \citep{olist2018dataset}. Olist aggregates
small sellers onto Brazil's large retail platforms and handles fulfilment, so each
order links a customer, one or more sellers, a basket, a payment, a delivery and
usually a review -- a structure supporting all three question families the brief
identifies: operational performance, customer experience and commercial outcomes.
The review below covers what Phase 2 depends on: fulfilment and satisfaction,
model families for tabular data, feature engineering on a relational order model,
and metrics that survive severe imbalance.

\subsection{Fulfilment, delivery promises and customer satisfaction}

The e-commerce service-quality literature is unusually direct about which lever
matters. \citet{parasuraman2005esqual} decompose electronic service quality into
efficiency, \emph{fulfilment}, system availability and privacy, and find
fulfilment -- the accuracy of delivery promises and the condition in which goods
arrive -- among the strongest contributors to perceived quality.
\citet{rao2011elsq} sharpen this into logistics terms across 260 online retailers.
Transactional evidence at platform scale now confirms it: across more than 15
million marketplace orders, \citet{deshpande2023logistics} trace a chain rather
than a single link -- logistics performance drives the service rating a customer
leaves, and that rating in turn moves seller sales. That chain makes a low-review
model a commercial instrument rather than a diagnostic.

The promise itself is a second, less obvious decision variable.
\citet{cui2024promise} show that committing to a faster delivery date raises sales
but also increases returns and depresses retention, so the optimal promise trades
demand against the cost of breaking it. An estimated delivery date is therefore a
business decision, not a forecast, and should not be expected to behave like one
(Section~\ref{sec:candidates}). Prediction likewise pays only when it feeds a
decision: \citet{kandula2021prescriptive} route machine-learned delivery outcomes
into a scheduling optimiser and report cost savings of up to 10.2\%.

Delivery is thus the dominant \emph{observable} input to how a customer rates a
transaction, which is why our second candidate problem is framed after delivery
rather than at checkout. The same literature warns what is \emph{not} observable
here -- product expectations, seller communication, returns handling -- none of
which appears in an order table, and any performance ceiling we report must be
read against that omission.

\subsection{Model families for tabular transactional data}

For structured data of this size the practical model space is well characterised.
Regularised linear and logistic models are the natural baseline: cheap,
interpretable, and -- as Section~\ref{sec:candidates} shows -- revealing about
distribution shift. Tree ensembles are the usual step up, whether by bagging
decorrelated trees \citep{breiman2001rf} or fitting residuals stagewise with
shrinkage \citep{friedman2001gbm, chen2016xgboost}. \citet{grinsztajn2022tabular}
benchmark both against modern deep architectures across 45 datasets and find
tree-based methods remain state of the art at this scale, the gap owing to
inductive biases better matched to irregular tabular targets.

We therefore adopt a small model budget: one linear family (ridge and
$\ell_2$-regularised logistic regression) and one tree ensemble (histogram-based
gradient boosting), reused across both framings. Section~\ref{sec:candidates}
justifies that complexity in both directions -- the ensemble earns its place on
lead time, while on repeat purchase neither family beats the baseline.

\subsection{Feature engineering on a relational order model}

Deriving features from nine tables raises a problem a flat file does not. Three of
them -- items, payments and reviews -- are one-to-many against orders, so a naive
join multiplies rows and reweights every aggregate towards multi-item orders. The
remedy is to aggregate each child relation to the analysis grain before joining.
Beyond that the usual transactional vocabulary applies: basket counts and sums,
ratios that normalise for order size, calendar encodings, and geographic distance
from postcode centroids. High-cardinality categoricals such as the 74 product
categories suit target encoding \citep{miccibarreca2001target} better than
one-hot, provided it is fitted inside the cross-validation loop.

Leakage is the dominant risk. \citet{kaufman2012leakage} formalise it as
information about the target unavailable at prediction time, and identify the two
forms that threaten this dataset: post-outcome features and improper partitioning.
\citet{kapoor2023leakage} show how routinely both survive peer review, cataloguing
eight types of leakage across 294 affected papers in 17 fields; ours are their
\emph{illegitimate feature} and \emph{no clean separation between training and
test data} categories, and we guard each with an explicit assertion rather than a
convention. Critically, this framing makes leakage relative to a \emph{prediction
moment} -- the principle our feature policy is built on, since delivery outcomes
are forbidden when forecasting delivery but admissible when predicting a review
not yet written. On partitioning, \citet{bergmeir2012cv} show that standard
$k$-fold cross-validation is invalid for temporally ordered data.

\subsection{Evaluation under imbalance and drift}

Both classification targets here are heavily imbalanced, and the brief is explicit
that accuracy alone misleads. \citet{saito2015prplot} show that precision--recall
curves are more informative than ROC curves under strong imbalance, because
ROC-AUC is insensitive to a large true-negative mass. We report PR-AUC against its
no-skill floor alongside balanced accuracy, treating raw accuracy as a diagnostic
rather than a headline.

The obvious next step -- correcting the imbalance -- is one we deliberately do not
take. \citet{goorbergh2022imbalance} test random under-sampling, over-sampling and
SMOTE \citep{chawla2002smote} and find none improves discrimination, while all
three badly miscalibrate the resulting probabilities, systematically
overestimating minority-class risk; they recommend adjusting the decision
threshold on uncorrected data instead. We therefore do not resample. We do weight
the loss by class, which serves the same end as their threshold shift, but it
carries the same caveat: we treat model outputs as \emph{decision scores rather
than calibrated risks}. Should Phase 2 need calibrated risk -- to rank at-risk
orders by expected cost -- we will fit on the observed prior and tune the
threshold, reporting calibration explicitly.

Regression brings a second complication. \citet{hyndman2006accuracy} argue for
judging errors against a naive benchmark, decisive here because delivery
performance improves dramatically across the window (Section~\ref{sec:temporal}).
Under such a level shift, $R^2$ against a baseline fitted on the past turns
negative even for a useful model, whereas mean absolute error stays interpretable.
This is textbook concept drift \citep{gama2014drift}, whose recommendation that
non-stationary models be refitted with respect to time shapes our Phase 2 plan.
All modelling uses \code{scikit-learn} pipelines \citep{pedregosa2011sklearn},
keeping preprocessing inside each fold.

\subsection{Conclusion: trends, gaps and opportunities}

Three themes emerge. The service-quality literature converges on \emph{fulfilment}
as the dominant observable driver of e-commerce satisfaction, and now traces it
through to seller revenue. The tabular machine-learning literature converges on a
small set of model families, with tree ensembles the reasonable ceiling at this
scale. And the methodological literature is emphatic -- increasingly so -- that on
imbalanced, temporally ordered data the choice of split, metric and imbalance
handling decides whether a number means anything at all.

The gap we address is narrow but concrete. Most work on this marketplace is
descriptive, or predicts review scores from features that would not exist when the
prediction is needed. Little of it is explicit about \emph{when} a model would run
-- a choice that changes both what is admissible and what is achievable, by a
factor of two in our own results. Making the prediction moment a first-class
design decision, and measuring its cost, is the contribution of
Section~\ref{sec:candidates}.
